\section{Covariance observations and the contradiction}
\label{sec:observations}

The matrices constructed in Section~\ref{sec:matrices} describe fluctuations
whose information content is incompatible with the logarithmic Sobolev
inequality of the stationary law.  We compare two bounds for the same
Gaussian covariance observation.  The first follows from logarithmic
Sobolev and is universal after integration over its scale.  The second
uses the frozen matrices to recover a scalar lognormal covariance
experiment.  Integration over scale is essential: it assigns to every
spectral block exactly the weight used by the tracial limit, regardless
of the sizes of its eigenvalues.

Throughout this section, limits refer to the counterexample sequence.
All observation times, caps, compact parameter intervals, and polynomial
degrees are fixed before taking that limit.  We suppress its index.
Recall that $Y_t$ is stationary with law $\eta$, that
$R(\eta)\le\lambda^{-1}$ with $\lambda\to1$, and that $R$ without a
measure argument denotes the positive matrix field.  In particular,
Equations~\eqref{eq:matrix-normalization}
and~\eqref{eq:matrix-saturation} give
\begin{equation}\label{eq:obs-energy-input}
 \E\tr R^2\longrightarrow1,
 \qquad \E|DR|^2\longrightarrow1.
\end{equation}
To distinguish the original Sylvester matrices from the subsequently
modified block matrices, denote the former by $\widehat H_i$.  Equations
\eqref{eq:matrix-sylvester} and~\eqref{eq:matrix-saturation} assert
\begin{equation}\label{eq:obs-sylvester-input}
 D_iR=\tfrac12(\widehat H_iR+R\widehat H_i),
 \qquad
 \sum_i\E\|[\widehat H_i,R]\|_{\HS}^2=o(1).
\end{equation}
Only the upper bound below uses these original matrices.

Fix $L\ge1$ and define
\begin{equation}\label{eq:obs-covariance-definition}
 f_L(q)=\frac{q}{1+q/L},\qquad
 C_T(y)=\Id+f_L(TR(y)^2),\qquad
 Z_T=C_T(Y_t)^{1/2}G,\quad T>0,
\end{equation}
where $G$ is an independent standard Gaussian vector.  Thus
$\Id\le C_T\le(1+L)\Id$.  We write $H(P\mid Q)$ for relative entropy
and $\MI(U;V)$ for mutual information, including the usual averaged
conditional version when a conditioning variable is displayed.

\subsection{The integrated information upper bound}

\begin{lemma}[Integrated spectral multipliers]
\label{lem:obs-multipliers}
For $v,w\ge0$, not both zero, put
\begin{equation}\label{eq:obs-multiplier}
 M_L(v,w)=(v+w)^2\int_0^\infty
 \frac{dT}{
 (1+Tv^2/L)(1+Tw^2/L)
 (1+(1+1/L)Tv^2)(1+(1+1/L)Tw^2)}.
\end{equation}
Set $M_L(0,0)=0$.  Then
\begin{equation}\label{eq:obs-multiplier-bounds}
 M_L(v,w)\le4\log(1+L)
 \quad\hbox{for all }v,w\ge0,
 \qquad
 M_L(v,w)\le9
 \quad\hbox{if }\tfrac12\le v/w\le2.
\end{equation}
The second bound concerns positive $v,w$ and is independent of $L$.
\end{lemma}

\begin{proof}
Write $a=1/L$ and $b=1+1/L$, so $b-a=1$.  If
$r=\max(v,w)>0$, discard the two denominator factors associated with
the smaller eigenvalue and substitute $q=Tr^2$.  This gives
\[
 M_L(v,w)\le4\int_0^\infty
 \frac{dq}{(1+aq)(1+bq)}
 =4\log(b/a)=4\log(1+L).
\]
For the identity in the middle, integration up to $Q$ gives
$\log(1+bQ)-\log(1+aQ)$, which tends to $\log(b/a)$.
If $v,w$ are positive and comparable as in the statement, let
$r_0=\min(v,w)$.  Discard the two factors with coefficient $a$, and
use $b\ge1$.  The remaining integral is at most
$\int_0^\infty(1+Tr_0^2)^{-2}\,dT=r_0^{-2}$.
Since $v+w\le3r_0$, this proves the second bound.
\end{proof}

\begin{proposition}[Universal integrated upper bound]
\label{prop:information-upper}
There is an absolute constant $C_0$ such that, for every fixed $t>0$
and $L\ge1$, the observations in
Equation~\eqref{eq:obs-covariance-definition} satisfy
\begin{equation}\label{eq:obs-information-upper}
 \limsup\int_0^\infty\MI(Y_t;Z_T)\,\frac{dT}{T^2}\le C_0.
\end{equation}
In particular, one may take $C_0=9/4$.
\end{proposition}

\begin{proof}
Let $p_T(z\mid y)$ denote the density of the centered Gaussian with
covariance $C_T(y)$, and let $p_T(z)$ be its marginal density.
For each output $z$, the posterior density with respect to $\eta$ is
$q_z(y)=p_T(z\mid y)/p_T(z)$.  Applying $2\lambda H\le\Ical$ to
this density and then averaging in $z$ gives
\begin{equation}\label{eq:obs-likelihood-lsi}
 2\lambda\MI(Y_t;Z_T)
 \le\E\big|D_y\log p_T(Z_T\mid Y_t)\big|^2.
\end{equation}
Here and below the expectation on the right includes the observation
noise.  The posterior normalization does not depend on $y$ and therefore
does not contribute to the derivative.

We compute that derivative explicitly.  For any positive covariance
$C$ and its derivative $D_iC$, differentiation of its determinant and
inverse yields
\[
 D_i\log p(z\mid y)
 =\tfrac12\left(
 z^*C^{-1}(D_iC)C^{-1}z-\tr(C^{-1}D_iC)\right).
\]
Conditional on $y$, set $B_i=C^{-1/2}(D_iC)C^{-1/2}$ and
$g=C^{-1/2}z$.  The vector $g$ is standard Gaussian.  Expanding its
fourth moments gives
$\E(g^*B_ig-\tr B_i)^2=2\tr B_i^2$; indeed the three Gaussian
pairings give $(\tr B_i)^2+2\tr B_i^2$ before subtracting the mean.
Consequently Equation~\eqref{eq:obs-likelihood-lsi} becomes
\begin{equation}\label{eq:obs-gaussian-fisher}
 \MI(Y_t;Z_T)\le\frac1{4\lambda}
 \E\sum_i\tr\big(C_T^{-1}D_iC_T\,C_T^{-1}D_iC_T\big).
\end{equation}

These calculations also apply to the Sobolev matrix field $R$.  To see
this without assuming additional regularity, fix $T,L$ and the finite
dimension of one member of the sequence.  The map
$R\mapsto C_T$ has bounded derivative on positive semidefinite
matrices, and its eigenvalues lie in $[1,1+L]$.  For each fixed $z$,
the map $C\mapsto p(z\mid C)^{1/2}$ has bounded derivative on this
covariance interval.  The Sobolev chain rule therefore puts
$q_z^{1/2}$ in the form domain; $p_T(z)>0$ and is a harmless fixed
normalization.  The Sobolev approximation at the opening of
Section~\ref{sec:calculus} extends the logarithmic Sobolev inequality
to this form-domain function.  The squared derivatives
integrate in $z$ to the nonnegative finite expression in
Equation~\eqref{eq:obs-gaussian-fisher}.  Approximation in the form
domain and Tonelli thus justify both that equation and the subsequent
scale integration.  The finite-dimensional bounds used for form-domain
membership do not enter its numerical coefficient.

In an eigenbasis of $R$, write $v,w$ for the row and column
eigenvalues.  The covariance function is
\[
 c_T(v)=1+f_L(Tv^2)
       =\frac{1+bTv^2}{1+aTv^2},
 \qquad a=1/L,\quad b=1+1/L.
\]
Its divided difference, including its derivative at $v=w$, is
\[
 \frac{c_T(v)-c_T(w)}{v-w}
 =\frac{T(v+w)}{(1+aTv^2)(1+aTw^2)}.
\]
Multiplying its square by $c_T(v)^{-1}c_T(w)^{-1}$ and integrating
against $dT/T^2$ gives exactly the multiplier in
Equation~\eqref{eq:obs-multiplier} for the squared entry
$|(D_iR)_{vw}|^2$.  When $v=w=0$, the covariance derivative is zero,
which explains the value assigned to $M_L(0,0)$.

It remains to obtain a bound independent of the fixed cap.  For
noncomparable pairs with at least one positive eigenvalue,
Equation~\eqref{eq:obs-sylvester-input} gives entrywise
\[
 |(D_iR)_{vw}|^2
 =\frac{(v+w)^2}{4(v-w)^2}
       |[\widehat H_i,R]_{vw}|^2
 \le\frac94|[\widehat H_i,R]_{vw}|^2.
\]
This includes a pair with exactly one zero eigenvalue.  Their total
expected derivative energy is therefore $o(1)$.  The kernel--kernel
derivative is zero almost everywhere, as established in the Sylvester
construction, and in any case has zero covariance multiplier.
Lemma~\ref{lem:obs-multipliers} now yields
\[
 \int_0^\infty\MI(Y_t;Z_T)\,\frac{dT}{T^2}
 \le\frac1{4\lambda}
 \left(9\E|DR|^2+4\log(1+L)\,o(1)\right).
\]
Take the sequence limit at this fixed $L$, and use
Equation~\eqref{eq:obs-energy-input}.  Stationarity makes the bound
independent of $t$.  This proves the proposition.  In particular, the
cap-dependent bound on noncomparable pairs is used only before its
vanishing error has been removed.
\end{proof}

\subsection{Stability on the entire scale interval}

We next transfer observations to the frozen covariance of
Proposition~\ref{prop:frozen-law}.  A bound on a compact interval of
$T$ would not suffice, because the useful scale of a block depends on
its eigenvalue.  The following comparison holds on all positive
scales.

\begin{samepage}
\begin{lemma}[Integrated covariance stability]
\label{lem:obs-integrated-stability}
For $T>0$ let $h_T(v)=\sqrt{1+f_L(Tv^2)}$, $v\ge0$, and put
\[
 K_L=\int_0^\infty
 \frac{dq}{(1+q/L)^3(1+(1+1/L)q)}.
\]
Then $0<K_L\le\log(1+L)$, and for arbitrary positive semidefinite
matrices $A,B$ of the same size,
\begin{equation}\label{eq:obs-matrix-stability}
 \int_0^\infty\|h_T(A)-h_T(B)\|_{\HS}^2\,\frac{dT}{T^2}
 \le K_L\|A-B\|_{\HS}^2.
\end{equation}
Let $\sigma$ be the frozen covariance in
Proposition~\ref{prop:frozen-law}, and couple
\[
 \widetilde Z_T=(\Id+f_L(T\sigma))^{1/2}G
\]
to $Z_T$ using the same independent Gaussian.  For every fixed $t,L$,
\begin{equation}\label{eq:obs-observation-stability}
 \int_0^\infty\E|Z_T-\widetilde Z_T|^2\,\frac{dT}{T^2}=o(1).
\end{equation}
\end{lemma}
\end{samepage}

\begin{proof}
Direct differentiation gives
\[
 h_T'(v)=\frac{Tv}
 {(1+Tv^2/L)^{3/2}(1+(1+1/L)Tv^2)^{1/2}}.
\]
For $v>0$, substitution $q=Tv^2$ therefore shows
\[
 \int_0^\infty|h_T'(v)|^2\,\frac{dT}{T^2}=K_L.
\]
Discarding two of the three factors $1+q/L$ bounds $K_L$ by the
integral evaluated in Lemma~\ref{lem:obs-multipliers}, namely
$\log(1+L)$.  If $v\ge w\ge0$, the fundamental theorem of calculus
and Cauchy--Schwarz give
\[
 |h_T(v)-h_T(w)|^2
 \le(v-w)\int_w^v|h_T'(r)|^2\,dr.
\]
Integrating in $T$ and using Tonelli proves the scalar bound
$K_L(v-w)^2$.  The value of the derivative integral at the single
endpoint $r=0$ does not affect this argument, so it covers $w=0$.

Choose an orthonormal eigenbasis $(e_a)$ of $A$ with eigenvalues $v_a$
and an orthonormal eigenbasis $(f_b)$ of $B$ with eigenvalues $w_b$.
The entry of their functional-calculus difference between these two
bases is
\[
 \langle e_a,(h_T(A)-h_T(B))f_b\rangle
   =(h_T(v_a)-h_T(w_b))\langle e_a,f_b\rangle.
\]
Multiply the scalar bound by $|\langle e_a,f_b\rangle|^2$ and sum
over $a,b$.  The sum on the right is
$K_L\|A-B\|_{\HS}^2$, proving
Equation~\eqref{eq:obs-matrix-stability}.  Finally, conditional on all
data except $G$, Gaussian isometry gives
\[
 \E_G|Z_T-\widetilde Z_T|^2
 =\|h_T(R(Y_t))-h_T(\sigma^{1/2})\|_{\HS}^2.
\]
Apply the matrix bound and the squared-mean convergence in
Proposition~\ref{prop:frozen-law}.  The resulting factor $K_L$ is
fixed while the sequence limit is taken.
\end{proof}

\subsection{Conditional information and blockwise tests}

Let $\mathcal D$ be the initial position, random spectral partition,
chosen orthonormal frame, and frozen block matrices from
Proposition~\ref{prop:frozen-law}.  Block formulas below use components
in this $\mathcal D$-measurable frame.  Every information term
$\MI(Y_t;Z_T)$ continues to refer to the original observation in
ambient coordinates; conditional on $\mathcal D$, conversion to
the block coordinates is invertible.  Orthogonal invariance makes
the components of the fresh Gaussian noise conditionally standard
and independent of the diffusion data.  Conditional on $\mathcal D$,
$N=B_t/\sqrt t$ is standard Gaussian.  On a positive block with
rounded eigenvalue $\beta$, write $H[v]=\sum_{i\in\beta}v_iH_i$.
The frozen covariance is
\begin{equation}\label{eq:obs-frozen-covariance}
 \sigma_\beta=\beta^2
       \exp\big(\sqrt{8t}\,H[N]-4t\Id\big).
\end{equation}
All $H_i$ here are the modified matrices supported on their initial
blocks.  Their bounds, from
Equations~\eqref{eq:flattening-bound} and~\eqref{eq:matrix-tracial-defects},
will be used in the form
\begin{align}
 &S=\sum_iH_i^2\le C\Id,
 \qquad \|H[v]\|_{\HS}^2\le C_1|v|^2,
 \label{eq:obs-block-bounds}\\
 &\tau(\Id)\longrightarrow1,
 \qquad\tau((S-\Id)^2)=o(1),
 \qquad\sum_{i,j}\tau(|[H_i,H_j]|^2)=o(1).
 \label{eq:obs-block-defects}
\end{align}
The constants $C,C_1\ge1$ are absolute.  The entries of the three-index
array $(H_i)_{ab}$ are also asymptotically symmetric in all three
indices in the weighted squared norm.  As before,
\[
 \tau(U)=\E\sum_\beta\beta^2\tr_\beta U_\beta,
\]
where expectation includes any Gaussian variables appearing in $U$.
No trace in this definition or in
Equation~\eqref{eq:obs-block-bounds} is normalized by block dimension.

\begin{lemma}[Conditional Ornstein--Uhlenbeck score]
\label{lem:obs-ou-score}
Let $N$ be a standard Gaussian vector conditional on $\mathcal D$,
and let $Z$ be any observation for which
$\MI(N;Z\mid\mathcal D)<\infty$.  With independent standard
Gaussian $G'$ set
\[
 N_u=\alpha N+\sqrt{1-\alpha^2}\,G',\qquad \alpha=e^{-u},\quad u>0.
\]
The relative score of its posterior, with respect to the standard
Gaussian marginal, is
\begin{equation}\label{eq:obs-ou-score}
 D\log\frac{p_{N_u\mid Z,\mathcal D}}{\gamma}(N_u)
 =\frac{\alpha}{1-\alpha^2}
   \E[N-\alpha N_u\mid N_u,Z,\mathcal D].
\end{equation}
For every compact interval $\mathcal J\subset(0,\infty)$,
\begin{equation}\label{eq:obs-ou-information}
 \MI(N;Z\mid\mathcal D)
 \ge\int_{\mathcal J}\frac{\alpha^2}{(1-\alpha^2)^2}
 \E\left|\E[N-\alpha N_u\mid N_u,Z,\mathcal D]\right|^2\,du.
\end{equation}
For the covariance observation in
Equation~\eqref{eq:obs-covariance-definition},
\begin{equation}\label{eq:obs-data-processing}
 \MI(Y_t;Z_T)\ge\MI(N;Z_T\mid\mathcal D).
\end{equation}
\end{lemma}

\begin{proof}
Put $q=1-\alpha^2$.  Given a posterior law $\pi(dn)$ of $N$ at
fixed $(Z,\mathcal D)$, the density at $y$ after adding the OU noise
is the integral of
$(2\pi q)^{-d/2}\exp(-|y-\alpha n|^2/(2q))$ against $\pi$.
Differentiating this kernel and dividing by its integral gives
\[
 D\log p(y\mid Z,\mathcal D)
 =-\frac{y-\alpha\E[N\mid N_u=y,Z,\mathcal D]}q.
\]
The score of the standard Gaussian is $-y$.  Subtracting it proves
Equation~\eqref{eq:obs-ou-score}.  Kernel differentiation is justified
first with $N$ truncated; on compact positive time intervals its
derivatives are bounded by Gaussian kernels times polynomials, and
truncation can then be removed by conditional second moments.

The OU semigroup has generator $\Delta-y\cdot D$.  For a positive
posterior density $r_u$ relative to $\gamma$, integration by parts
gives
\[
 \frac{d}{du}\int r_u\log r_u\,d\gamma
 =-\int\frac{|Dr_u|^2}{r_u}\,d\gamma.
\]
This is the entropy dissipation calculation of
Lemma~\ref{lem:entropy-calculus}; positive-time smoothing, truncation
of the density and spatial cutoffs justify it for finite-entropy
posteriors.  Averaging over $(Z,\mathcal D)$ and integrating over
$\mathcal J$ bounds its nonnegative right side by the initial
conditional information.  Substitution of the score formula proves
Equation~\eqref{eq:obs-ou-information}.

For the last assertion, observation noise is independent of the
initial data and diffusion path, and its law depends on them only
through $Y_t$.  Hence
$(\mathcal D,N)\longrightarrow Y_t\longrightarrow Z_T$ is a Markov
chain.  The chain rule and data processing give
\[
 \MI(N;Z_T\mid\mathcal D)
 \le\MI(\mathcal D,N;Z_T)\le\MI(Y_t;Z_T).
\]
These informations are finite for each member of the sequence:
the covariance interval $[\Id,(1+L)\Id]$ bounds the output entropy
from above by that of a Gaussian of covariance $(1+L)\Id$ and
bounds its conditional entropy from below by that of a standard
Gaussian.  Thus $\MI(Y_t;Z_T)\le(d/2)\log(1+L)$ in dimension $d$.
This finiteness assertion is used only to justify the identities.
\end{proof}

\begin{lemma}[Block tests and the scale Jacobian]
\label{lem:obs-block-tests}
Fix compact intervals $\mathcal J,\mathcal X\subset(0,\infty)$, and
write $\kappa(u)=\alpha^2/(1-\alpha^2)^2$, where $\alpha=e^{-u}$.
Let $P(u,x,a,b)$ be a real polynomial, evaluated on matrices with any
one fixed ordering of each monomial.  On each positive block define
\[
 A_\beta=H_\beta[N_\beta],\qquad
 A_{u,\beta}=H_\beta[(N_u)_\beta],\qquad
 B_{\beta,x}=H_\beta[(Z_{x/\beta^2})_\beta],
\]
and $P_{\beta,u,x}=P(u,x,A_{u,\beta},B_{\beta,x})$.  Then
\begin{align}
 \int_0^\infty\MI(Y_t;Z_T)\,\frac{dT}{T^2}
 \ge{}&\int_{\mathcal J}\int_{\mathcal X}\kappa(u)
 \E\sum_\beta\beta^2\tr_\beta
 \bigl(2(A_\beta-\alpha A_{u,\beta})P_{\beta,u,x}
 \notag\\[-2pt]
 &\hspace{43mm}-C_1P_{\beta,u,x}^*P_{\beta,u,x}\bigr)
 \,\frac{dx}{x^2}\,du.
 \label{eq:obs-block-test-bound}
\end{align}
Here the constant $C_1$ is the absolute constant in
Equation~\eqref{eq:obs-block-bounds}.
\end{lemma}

\begin{proof}
For a block, consider the linear map
$\mathcal H:v\mapsto H[v]$ between its Euclidean coordinate space
and its matrix space with the ordinary Hilbert--Schmidt norm.
Its adjoint is
\[
 (\mathcal H^*Q)_i=\tr(H_iQ),
\]
because $H_i$ is real symmetric.  Taking the adjoint in
Equation~\eqref{eq:obs-block-bounds} gives
\begin{equation}\label{eq:obs-adjoint-bound}
 \sum_i|\tr(H_iQ)|^2\le C_1\tr(Q^*Q).
\end{equation}
This step has no factor involving the dimension of the block.

At fixed $u,T$, select precisely those positive blocks for which
$x=T\beta^2\in\mathcal X$.  On a selected block set
\[
 V_i=\tr_\beta
 \bigl(H_iP(u,T\beta^2,H[(N_u)_\beta],H[(Z_T)_\beta])\bigr),
 \qquad i\in\beta,
\]
and set all other components, including those in the kernel, equal
to zero.  This vector is measurable with respect to
$(N_u,Z_T,\mathcal D)$.  If
$m=\E[N-\alpha N_u\mid N_u,Z_T,\mathcal D]$, completing the square
gives
\[
 \E|m|^2\ge2\E(N-\alpha N_u)\cdot V-\E|V|^2.
\]
The coordinates of distinct selected blocks are orthogonal.  Their
costs therefore add, and Equation~\eqref{eq:obs-adjoint-bound} bounds
the cost by $C_1$ times the sum of their unnormalized squared matrix
norms.  Their linear terms contract exactly to
\[
 (N-\alpha N_u)\cdot V
 =\sum_{\beta\text{ selected}}\tr_\beta
 \bigl((H[N_\beta]-\alpha H[(N_u)_\beta])P\bigr).
\]
Substitute these two statements in
Lemma~\ref{lem:obs-ou-score} and integrate in $T$.

For each realized positive block separately, the change of variable
$x=T\beta^2$ has the exact Jacobian
\begin{equation}\label{eq:obs-scale-jacobian}
 T=\frac{x}{\beta^2},\qquad dT=\frac{dx}{\beta^2},\qquad
 \boxed{\displaystyle\frac{dT}{T^2}=\beta^2\frac{dx}{x^2}}.
\end{equation}
Thus the block selection becomes $x\in\mathcal X$, and the
coefficient of each block trace becomes precisely $\beta^2$.
This proves Equation~\eqref{eq:obs-block-test-bound}.

For completeness, all its polynomial terms are integrable on the
compact parameter rectangle.  The moment bounds of
Lemma~\ref{lem:matrix-gaussian-moments} control $A_\beta$ and
$A_{u,\beta}$ under the weighted trace.  Conditional on the
initial data, diffusion path, and OU noise, $B_{\beta,x}$ is a
Gaussian matrix series whose vector covariance is at most
$(1+L)\Id$.  If its coefficients are $K_j=\sum_i Q_{ij}H_i$
for a square root $Q$ of that covariance, then for every vector $v$
\[
 \sum_j|K_jv|^2\le(1+L)\sum_i|H_iv|^2\le C(1+L)|v|^2.
\]
The same Gaussian even-moment proof therefore gives, for every fixed
integer $p\ge1$,
\begin{equation}\label{eq:obs-gaussian-moment-bound}
 \E_G\tr_\beta|B_{\beta,x}|^{2p}
 \le[C(1+L)]^p(2p-1)!!\tr_\beta\Id.
\end{equation}
Trace H\"older bounds each mixed word by its factors' moments.
After multiplying by $\beta^2$ and summing, the right side is a
constant times $\tau(\Id)$, which is finite and bounded along the
sequence.  The measure $\kappa(u)\,du\,dx/x^2$ is finite on the
chosen rectangle.  Tonelli applies to the square terms, while these
moment bounds justify Fubini for the linear terms.  One can
equivalently first truncate $V$ in the conditional-mean inequality
and then remove that truncation in this integrated squared norm.
\end{proof}

\subsection{Transfer of the polynomial tests to scalar observations}

The scale substitution has now produced the tracial weights of the
matrix limit.  We next identify that limit for the matrices containing
the observation.  Two replacements are needed: first replace $Z_T$
by its frozen version on the whole scale interval, and then use the
approximate symmetry of the three-index array to move a bounded
matrix function through $H[\,\cdot\,]$.

\begin{lemma}[Replacement in integrated polynomial tests]
\label{lem:obs-polynomial-replacement}
In the notation of Lemma~\ref{lem:obs-block-tests}, put
$\widetilde B_{\beta,x}
=H_\beta[(\widetilde Z_{x/\beta^2})_\beta]$.
For every fixed polynomial in
$u,x,A_\beta,A_{u,\beta},B_{\beta,x}$ and their adjoints, its
weighted trace integral over $\mathcal J\times\mathcal X$ changes
by $o(1)$ if $B_{\beta,x}$ is replaced by
$\widetilde B_{\beta,x}$.
\end{lemma}

\begin{proof}
The flattening bound and the inverse substitution in
Equation~\eqref{eq:obs-scale-jacobian} give
\begin{align*}
 &\E\sum_\beta\beta^2\int_{\mathcal X}
 \|B_{\beta,x}-\widetilde B_{\beta,x}\|_{\HS}^2\,\frac{dx}{x^2}
 \\
 &\quad=\int_0^\infty\E\sum_{\beta:T\beta^2\in\mathcal X}
 \|H_\beta[(Z_T-\widetilde Z_T)_\beta]\|_{\HS}^2\,\frac{dT}{T^2}
 \\
 &\quad\le C_1\int_0^\infty
 \E|Z_T-\widetilde Z_T|^2\,\frac{dT}{T^2}=o(1),
\end{align*}
by Lemma~\ref{lem:obs-integrated-stability}.  All projected errors
are summed before this comparison; no small block weight is divided
out.  Multiplication by the bounded weight $\kappa(u)$ and
integration in the compact interval $\mathcal J$ preserve $o(1)$.

To pass from this squared norm to arbitrary fixed products, include
the parameter integrals in the finite tracial functional
\[
 \mathcal T(U)=\int_{\mathcal J}\int_{\mathcal X}
 \kappa(u)\E\sum_\beta\beta^2\tr_\beta U_{\beta,u,x}
 \,\frac{dx}{x^2}\,du.
\]
Every matrix in the statement has bounded $L^p(\mathcal T)$ norm
for every fixed finite $p$.  For $B$ this was proved in
Equation~\eqref{eq:obs-gaussian-moment-bound}; for $\widetilde B$
the same proof applies because its covariance is also at most
$(1+L)\Id$.  The remaining matrices are covered by the Gaussian
moment lemma.  For any $q>p\ge2$, interpolation with
$1/p=\theta/2+(1-\theta)/q$ gives
\[
 \|B-\widetilde B\|_{L^p(\mathcal T)}
 \le\|B-\widetilde B\|_{L^2(\mathcal T)}^\theta
       \|B-\widetilde B\|_{L^q(\mathcal T)}^{1-\theta}=o(1).
\]
For lower powers, H\"older and boundedness of
$\mathcal T(\Id)$ give the same conclusion.  Expand the difference
of two words of length $r$ by replacing one factor at a time.
Trace H\"older with all exponents $r$ bounds each resulting trace
by the product of its factors' $L^r(\mathcal T)$ norms.  The changed
factor tends to zero and the others remain bounded.  There are
only finitely many such terms.  This proves the claim, including
polynomial squared norms and linear terms with $A-\alpha A_u$.
\end{proof}

For fixed $t,L$, define the scalar function
\begin{equation}\label{eq:obs-g-function}
 g_x(v)=\sqrt{1+f_L\big(xe^{\sqrt{8t}\,v-4t}\big)},
 \qquad x>0,\ v\in\R.
\end{equation}
Then $1\le g_x\le\sqrt{1+L}$.  By
Equation~\eqref{eq:obs-frozen-covariance}, the frozen observation on
a block at $T=x/\beta^2$ is exactly
$g_x(A_\beta)G_\beta$.  Thus
$\widetilde B_{\beta,x}=H_\beta[g_x(A_\beta)G_\beta]$.

\begin{lemma}[Polynomial density under exponential moments]
\label{lem:obs-polynomial-density}
Let $\nu$ be a finite positive measure on $\R^k$ with
$\int e^{a|w|}\,d\nu(w)<\infty$ for every $a>0$.  Polynomials
are dense in $L^p(\nu)$ for every $1\le p<\infty$.
\end{lemma}

\begin{proof}
Suppose a continuous linear functional on $L^p(\nu)$ annihilates
all polynomials.  Its representing density $F$ belongs to
$L^{p'}(\nu)$, where $p'=p/(p-1)$ and $p'=\infty$ if $p=1$.
H\"older and the hypothesis show that the finite signed measure
$F\,d\nu$ has exponential moments of every linear magnitude.
For each $\xi\in\R^k$, the function
\[
 z\longmapsto\int e^{z\xi\cdot w}F(w)\,d\nu(w)
\]
is entire: on every compact set of complex $z$, its derivatives
are dominated by an integrable exponential.  All its derivatives
at zero vanish because $(\xi\cdot w)^j$ is a polynomial.
Its Taylor series and hence the entire function are zero.
Evaluating at $z=\mathrm i$ proves that the Fourier transform of
$F\,d\nu$ vanishes.  Fourier uniqueness for finite measures can
here be seen by convolving with a centered Gaussian: the convolution
has zero Fourier transform and an integrable Fourier inversion
formula, so is zero for every positive Gaussian variance; letting
that variance tend to zero shows the original measure is zero.
Thus $F=0$ almost everywhere.  Duality, or equivalently separation
of a proper closed subspace, proves the density assertion.
\end{proof}

\begin{lemma}[Nonlinear mixed trace law]
\label{lem:obs-nonlinear-trace-law}
Let $n,g',g$ be independent scalar standard Gaussians and let
$n_u=\alpha n+\sqrt{1-\alpha^2}\,g'$.  For every fixed mixed
polynomial, the weighted trace moments of
\[
 \big(A_\beta,A_{u,\beta},\widetilde B_{\beta,x}\big)
\]
converge to the corresponding scalar moments of
\[
 \big(n,n_u,g_x(n)g\big).
\]
The assertion holds after integration over compact positive $u,x$
intervals against $\kappa(u)\,du\,dx/x^2$, and permits fixed
polynomial dependence on $u,x$.  Adjoints in a matrix word become
the identical real scalar factors in its limit.
\end{lemma}

\begin{proof}
First we prove the matrix comparison
\begin{equation}\label{eq:obs-index-swap-comparison}
 H[g_x(A)G]-g_x(A)H[G]\longrightarrow0
 \quad\hbox{in weighted squared Hilbert--Schmidt norm}.
\end{equation}
All notation in this calculation is on one block.  Define
\[
 J(v)_{ab}=(H_bv)_a,
 \qquad D(v)=H[v]-J(v),
 \qquad \Delta_{abi}=(H_i)_{ab}-(H_b)_{ai}.
\]
Put $q=g_x(A)$ and $B=H[G]$.  The exact columnwise identity is
\begin{equation}\label{eq:obs-index-swap-identity}
 H[qG]-qB=D(qG)+K-qD(G),
 \qquad K_{\cdot b}=[H_b,q]G.
\end{equation}
Indeed the $b$th column of $J(qG)-qJ(G)$ is
$H_bqG-qH_bG$.  Conditional Gaussian averaging therefore bounds
the square of the left side by
\begin{equation}\label{eq:obs-index-swap-error}
 6(1+L)\sum_{a,b,i}|\Delta_{abi}|^2
       +3\sum_b\|[H_b,g_x(A)]\|_{\HS}^2.
\end{equation}
For the first and last terms in
Equation~\eqref{eq:obs-index-swap-identity}, we used
$\|q\|_{\op}\le\sqrt{1+L}$ and the Gaussian isometry for the
linear map $D$.

The function $g_x$ is Lipschitz with a bound depending only on
$t,L$, uniformly for all $x>0$.  In fact, with
$a=\sqrt{8t}$ and $r=xe^{av-4t}/L$, differentiation gives
\[
 |g_x'(v)|
 =\frac{aLr}{2(1+r)^2g_x(v)}\le\frac{aL}{8}.
\]
In an eigenbasis of $A$ with eigenvalues $r_a$, the squared
commutator norm is
\[
 \|[H_b,g_x(A)]\|_{\HS}^2
 =\sum_{a,c}|(H_b)_{ac}|^2|g_x(r_c)-g_x(r_a)|^2
 \le\operatorname{Lip}(g_x)^2\|[H_b,A]\|_{\HS}^2.
\]
This is a scalar Lipschitz estimate for matrix entries in the
Hilbert--Schmidt norm.  Averaging $A=\sum_iN_iH_i$ over $N$ gives
\[
 \E_N\sum_b\|[H_b,A]\|_{\HS}^2
 =\sum_{b,i}\|[H_b,H_i]\|_{\HS}^2.
\]
Multiply Equation~\eqref{eq:obs-index-swap-error} by the block
weight, sum, and average.  The total symmetry defect and
Equation~\eqref{eq:obs-block-defects} prove
Equation~\eqref{eq:obs-index-swap-comparison}, uniformly in $x$ at
fixed $t,L$.

Every fixed higher moment of both sides of this comparison is
bounded.  Conditional on $A$ and the initial data,
$H[g_x(A)G]$ is a Gaussian matrix series with coefficient-square
sum at most $(1+L)S$, by the quadratic-form argument used in
Equation~\eqref{eq:obs-gaussian-moment-bound}.  Also the Schatten
ideal inequality gives
\[
 \|g_x(A)H[G]\|_{p}\le\sqrt{1+L}\,\|H[G]\|_{p}
\]
on each block.  Interpolation and telescoping products, exactly as
in Lemma~\ref{lem:obs-polynomial-replacement}, show that these
two matrices may be interchanged in every fixed polynomial trace.

It remains to identify traces containing $g_x(A)H[G]$.
Lemma~\ref{lem:matrix-gaussian-moments} gives the joint scalar
Gaussian moment law for $A=H[N]$, $H[G']$, and $H[G]$.
In particular, $A_u=\alpha A+\sqrt{1-\alpha^2}H[G']$ has the
required joint law with them.  To include the bounded continuous
function $g_x(A)$, fix $x$ and a mixed word of length $r$.
Choose a polynomial $p$ approximating $g_x$ in a sufficiently high
finite $L^q(\gamma_1)$ norm, where $\gamma_1$ is the scalar
standard Gaussian law.  Choose $q$ larger than all the H\"older
exponents needed for this fixed word, and take the approximation
error at most $1$; then $\|p\|_{L^q(\gamma_1)}
\le\sqrt{1+L}+1$.  Such approximations exist by the
polynomial-density argument in
Lemma~\ref{lem:obs-polynomial-density}, applied to one Gaussian
coordinate.  The scalar spectral measures of $A$ under $\tau$
converge to $\gamma_1$ by the moment and exponential-moment bounds
in Lemma~\ref{lem:matrix-gaussian-moments}.  Their higher moments
are uniformly bounded.  Consequently, for even $q$,
\begin{equation}\label{eq:obs-functional-approximation}
 \tau\bigl(|g_x(A)-p(A)|^q\bigr)
 \longrightarrow\E|g_x(n)-p(n)|^q.
\end{equation}
To justify this last passage directly, cut the scalar spectral
variable off at $|v|\le M$, use weak convergence for the bounded
continuous integrand there, and bound the complement by a higher
uniform moment of $A$.  Since $g_x$ is bounded and $p$ has fixed
degree, that tail tends to zero as $M\to\infty$.

Replace each occurrence of $g_x(A)$ by $p(A)$.  Telescoping the
word and applying trace H\"older bounds the replacement error
by a constant times the $L^q$ error in
Equation~\eqref{eq:obs-functional-approximation}, with all other
factors controlled by fixed higher moments.  For factors $p(A)$,
the preceding $L^q$ bound and moment convergence give a bound
independent of how accurately $p$ approximates $g_x$.
Thus the error constant stays bounded as that approximation is
improved.  The remaining word
is polynomial in the three independent Gaussian matrix series,
so its trace has the scalar Gaussian limit.  Let the scalar
approximation error tend to zero.  This proves the claimed limit
at each fixed $u,x$; the scalar third coordinate is
$g_x(n)g$.

Finally, the moment bounds for the original, bounded function
$g_x$ are uniform on the compact parameter rectangle, as are the
bounds on polynomial coefficients in $u,x$ and on $\alpha$.
They dominate each mixed trace by an integrable constant.
Dominated convergence thus permits parameter integration.  This
last argument uses the uniform bounds for $g_x$ itself, and does
not require choosing one polynomial approximant uniformly in $x$.
\end{proof}

The preceding lemmas identify the polynomial tests of the actual
observation with those of the scalar variance mixture.  We can now
optimize the tests and integrate the scalar OU dissipation.  The
constant in the resulting comparison is exactly the constant of
the unnormalized flattening estimate.

\begin{proposition}[Scalar information lower bound]
\label{prop:information-lower}
For every fixed $t>0$ and $L\ge1$, let $g_x$ be defined by
Equation~\eqref{eq:obs-g-function}, and let $n,g$ be independent
scalar standard Gaussians.  Then
\begin{equation}\label{eq:obs-information-lower}
 \liminf\int_0^\infty\MI(Y_t;Z_T)\,\frac{dT}{T^2}
 \ge C_1^{-1}\int_0^\infty
       \MI\big(n;g_x(n)g\big)\,\frac{dx}{x^2}.
\end{equation}
The constant $C_1$ is universal and is independent of $t,L$.
\end{proposition}

\begin{proof}
Fix compact positive intervals $\mathcal J,\mathcal X$, and let
$n,g',g$ be independent scalar standard Gaussians.  Write
$n_u=\alpha n+\sqrt{1-\alpha^2}\,g'$, and fix a scalar polynomial
$P(u,x,a,b)$.  Give its matrix version any fixed
ordering of the factors of each monomial.  Lemmas
\ref{lem:obs-block-tests}, \ref{lem:obs-polynomial-replacement},
and~\ref{lem:obs-nonlinear-trace-law}, in that order, imply
\begin{align}
 &\liminf\int_0^\infty\MI(Y_t;Z_T)\,\frac{dT}{T^2}
 \notag\\
 &\quad\ge\int_{\mathcal J}\int_{\mathcal X}\kappa(u)
 \E\big[2(n-\alpha n_u)P(u,x,n_u,g_x(n)g)
 \notag\\[-2pt]
 &\hspace{58mm}-C_1P(u,x,n_u,g_x(n)g)^2\big]
       \,\frac{dx}{x^2}\,du.
 \label{eq:obs-scalar-test-bound}
\end{align}
If the matrix polynomial is not symmetric, its prelimit cost is
$\tr(P^*P)$ as in Equation~\eqref{eq:obs-block-test-bound}.
The mixed trace law takes this cost to the real scalar $P^2$,
so symmetry of the chosen ordering is unnecessary.

The sequence limit in Equation~\eqref{eq:obs-scalar-test-bound}
has been taken for each fixed polynomial.  We now optimize the
tests in one squared-mean space containing the compact parameters.
Give $(u,x)$ the finite measure $\kappa(u)\,du\,dx/x^2$ on
$\mathcal J\times\mathcal X$, independently of the scalar
Gaussians $n,g',g$, and put
\begin{equation}\label{eq:obs-scalar-observed-vector}
 W=(u,x,n_u,g_x(n)g),\qquad r=n-\alpha n_u.
\end{equation}
The first two coordinates of $W$ are bounded.  At each fixed
$u,x$, the variable $n_u$ is standard Gaussian, and
$|g_x(n)g|\le\sqrt{1+L}|g|$.
Cauchy--Schwarz therefore gives all linear exponential moments
for its finite law, uniformly on the compact parameter rectangle.
Lemma~\ref{lem:obs-polynomial-density} shows that polynomials in
the four coordinates of $W$ are dense in this $L^2$ space.

Since $\E[r^2\mid u,x]=1-\alpha^2\le1$, the conditional mean
$m(W)=\E[r\mid W]$ belongs to the same space.  For the finite
joint measure just defined, completion of the square gives
\[
 \int\big(2rP(W)-C_1P(W)^2\big)
 =C_1^{-1}\int m^2-C_1\int(P-m/C_1)^2.
\]
Density makes the infimum of the final squared error zero, using
single four-variable polynomials in the integrated space.
Taking their supremum and then disintegrating with respect to
$(u,x)$ turns the right side of
Equation~\eqref{eq:obs-scalar-test-bound} into
\begin{equation}\label{eq:obs-scalar-score-lower}
 C_1^{-1}\int_{\mathcal J}\int_{\mathcal X}
 \frac{\alpha^2}{(1-\alpha^2)^2}
 \E\left|\E[n-\alpha n_u\mid n_u,g_x(n)g]\right|^2
 \,\frac{dx}{x^2}\,du.
\end{equation}
The conditional expectation here is taken at each fixed $u,x$.
All approximation constants used before this optimization may
depend on the fixed polynomial, $t,L$, and the compact intervals.
The coefficient $C_1$ depends only on
Equation~\eqref{eq:obs-adjoint-bound}.

We verify the scalar information identity including its endpoints.
Fix $x>0$ and write $z_x=g_x(n)g$.  Its information is finite:
conditional variance lies in $[1,1+L]$, so Gaussian entropy
comparison gives
\begin{equation}\label{eq:obs-scalar-information-finite}
 0\le\MI(n;z_x)\le\tfrac12\log(1+L).
\end{equation}
Indeed $\E z_x^2\le1+L$ bounds its entropy above by
$\frac12\log(2\pi e(1+L))$, while its conditional entropy is
$\frac12\log(2\pi e)+\E\log g_x(n)$ and is at least
$\frac12\log(2\pi e)$.
Apply the scalar version of Lemma~\ref{lem:obs-ou-score} and its
entropy dissipation identity between times $a$ and $b$, where
$0<a<b<\infty$.  This gives
\begin{align*}
 \MI(n_a;z_x)-\MI(n_b;z_x)
 =\int_a^b\frac{\alpha^2}{(1-\alpha^2)^2}
 \E\left|\E[n-\alpha n_u\mid n_u,z_x]\right|^2\,du.
\end{align*}
As $a\downarrow0$, the joint pair $(n_a,z_x)$ converges in law
to $(n,z_x)$.  Lower semicontinuity of relative entropy, applied
to the joint law and the product of its marginals, gives
$\MI(n;z_x)\le\liminf_{a\downarrow0}\MI(n_a;z_x)$.
The reverse inequality follows from the Markov chain
$z_x\longrightarrow n\longrightarrow n_a$ and data processing.
Thus the initial endpoint is exactly $\MI(n;z_x)$.
At the other endpoint, the same chain gives
\[
 0\le\MI(n_u;z_x)\le\MI(n_u;n)
       =-\tfrac12\log(1-\alpha^2)\longrightarrow0.
\]
The equality follows by subtracting the conditional Gaussian
entropy of variance $1-\alpha^2$ from that of the standard
Gaussian marginal.  We have proved
\begin{equation}\label{eq:obs-scalar-ou-identity}
 \MI(n;z_x)=\int_0^\infty
 \frac{\alpha^2}{(1-\alpha^2)^2}
 \E\left|\E[n-\alpha n_u\mid n_u,z_x]\right|^2\,du.
\end{equation}

Finally exhaust both positive parameter axes by compact intervals
in Equation~\eqref{eq:obs-scalar-score-lower}.  Its integrand is
nonnegative.  Monotone convergence and
Equation~\eqref{eq:obs-scalar-ou-identity} prove
Equation~\eqref{eq:obs-information-lower}.  In particular, the
polynomial optimization and the exhaustion take place after the
counterexample-sequence limit, so neither requires a uniform
approximation over unbounded parameters.
\end{proof}

\subsection{Divergence of the scalar scale integral}

The upper and lower comparisons are now in the same information
units, with universal constants $C_0,C_1$.  To contradict them, it
suffices to let the fixed diffusion time and the fixed covariance
cap grow together.  The following elementary test detects the
large values of the scalar lognormal variance.

\begin{lemma}[Scalar information divergence]
\label{lem:scalar-divergence}
Let $n,g$ be independent scalar standard Gaussians and, for
$t\ge1$, define
\[
 \ell=\exp(\sqrt{8t}\,n-4t),\qquad
 z_x=\sqrt{1+f_t(x\ell)}\,g.
\]
There are absolute constants $c>0$ and $h_0\ge1$ such that, for
all $t\ge h_0$,
\begin{equation}\label{eq:obs-scalar-divergence}
 \int_0^\infty\MI(n;z_x)\,\frac{dx}{x^2}
 \ge c\E\big[\ell\1_{\{\ell\ge e^{2t}\}}\big]
                  \log(t/h_0)
 \ge\frac c2\log(t/h_0).
\end{equation}
\end{lemma}

\begin{proof}
The Gaussian moment generating formula gives $\E\ell=1$.
Let $P_{x,n}$ be the conditional law of $z_x$ given $n$, and
let $Q_x$ be its marginal law.  Thus
\[
 \MI(n;z_x)=\E_n H(P_{x,n}\mid Q_x).
\]
Set $p_0=\Prob\{|g|\ge\sqrt2\}>0$.  Suppose first that
$\ell\ge e^{2t}$, and choose a scale with
$h=x\ell\in[h_0,t]$.  Since $h\le t$,
\[
 \Var(z_x\mid n)=1+\frac{h}{1+h/t}\ge h/2.
\]
It follows that the event $A_h=\{|z|\ge\sqrt h\}$ has
conditional probability $p=P_{x,n}(A_h)\ge p_0$.

To bound its marginal probability, use an independent variable
$\ell'$ with the law of $\ell$.  On $x\ell'\le1$, the
conditional variance of the mixture component is at most $2$,
so its probability of $A_h$ is at most $2e^{-h/4}$ by the
Gaussian exponential bound.  On the complementary event,
Markov's inequality and $\E\ell'=1$ give probability at most
$x$.  Hence
\begin{equation}\label{eq:obs-binary-marginal-bound}
 q=Q_x(A_h)\le2e^{-h/4}+x
       \le2e^{-h/4}+te^{-2t}\le3e^{-h/4}.
\end{equation}
The last inequality holds for $t\ge1$ and $h\le t$, because
$te^{-2t}\le e^{-t/4}\le e^{-h/4}$.

Apply data processing to the event indicator.  The binary
relative entropy is bounded below by
\begin{align*}
 H(P_{x,n}\mid Q_x)
 &\ge p\log(p/q)+(1-p)\log\big((1-p)/(1-q)\big)\\
 &\ge p\log(1/q)-\log2\\
 &\ge p_0(h/4-\log3)-\log2.
\end{align*}
Choose once and for all
\[
 h_0\ge\max\{1,\,8\log3+8\log2/p_0\},
 \qquad c=p_0/8.
\]
The preceding expression is then at least $ch$ for
$h\in[h_0,t]$.  This entropy estimate is made separately for
each conditional law $P_{x,n}$; the chosen event may therefore
depend on $n$ through $h=x\ell$.

Nonnegativity of relative entropy permits restricting the scale
integral to these choices of $n$ and $x$.  Tonelli and the change
of variable $x=h/\ell$ yield
\begin{align*}
 \int_0^\infty\MI(n;z_x)\,\frac{dx}{x^2}
 &\ge\E\left[\1_{\{\ell\ge e^{2t}\}}
       \int_{h_0/\ell}^{t/\ell}
           H(P_{x,n}\mid Q_x)\,\frac{dx}{x^2}\right]\\
 &\ge c\E\left[\ell\1_{\{\ell\ge e^{2t}\}}
                         \int_{h_0}^t\frac{dh}{h}\right]\\
 &=c\E\big[\ell\1_{\{\ell\ge e^{2t}\}}\big]\log(t/h_0).
\end{align*}
Under the probability change of measure with density $\ell$,
completion of the square makes $n$ a Gaussian of mean
$\sqrt{8t}$ and variance $1$.  Therefore $\log\ell$ under that
measure is Gaussian of mean $4t$ and variance $8t$.  If $\Phi$
is the standard Gaussian distribution function, then
\[
 \E\big[\ell\1_{\{\ell\ge e^{2t}\}}\big]
   =\Phi(\sqrt{t/2})\ge\tfrac12.
\]
This proves both inequalities in
Equation~\eqref{eq:obs-scalar-divergence}.
\end{proof}

\begin{proof}[Proof of Theorem~\ref{thm:main}]
Suppose that the asserted dimension-free bound fails.
Lemma~\ref{lem:critical-normalization} gives compact centered
counterexamples with normalized logarithmic Sobolev constant $1$
and linear subgaussian parameter tending to zero.
Proposition~\ref{prop:regularized-constants} shows that their
positive-time regularized constants converge to $1$.
Proposition~\ref{prop:critical-sequence} supplies the stationary
laws and near-eigenfunction fields at the fixed regularization
time $s=1/8$, in either the strict entropy branch or the spectral
gap branch.  Proposition~\ref{prop:prediction} gives the exact
prediction curve in both cases.  Proposition~\ref{prop:hierarchy}
then gives the symmetric tensor hierarchy, and
Proposition~\ref{prop:frozen-law} constructs its matrix square
root and frozen block description.  These are precisely the
inputs used in the present section.

For each fixed $t>0$ and $L\ge1$, combine
Propositions~\ref{prop:information-upper}
and~\ref{prop:information-lower}.  They give
\begin{equation}\label{eq:obs-final-scalar-bound}
 \int_0^\infty
 \MI\left(n;
 \sqrt{1+f_L(xe^{\sqrt{8t}\,n-4t})}\,g\right)
       \,\frac{dx}{x^2}\le C_0C_1,
\end{equation}
where $C_0C_1$ is independent of both fixed parameters.
Choose an arbitrarily large but finite $t\ge h_0$ and the
fixed cap $L=t$.  Lemma~\ref{lem:scalar-divergence} makes the
left side at least $(c/2)\log(t/h_0)$, which exceeds $C_0C_1$
for a sufficiently large choice of $t$.  This is a contradiction.

All counterexample limits used to obtain
Equation~\eqref{eq:obs-final-scalar-bound} were taken with this
$t,L$ fixed.  Only afterward was their numerical size chosen.
Thus no uniform estimate on the intermediate errors as
$t,L\to\infty$ is needed.  The contradiction excludes an
unbounded ratio $R(\mu)/a^2$ over the measures in the theorem.
Finally the convention
$\Ent_\mu(f^2)\le2R(\mu)\int|Df|^2\,d\mu$
absorbs the factor $2$ into a universal constant and gives the
asserted inequality for every smooth compactly supported real
test function.
\end{proof}
